\documentclass{article}

% Pass the round citation option to natbib BEFORE the NeurIPS style auto-loads it.
\PassOptionsToPackage{round}{natbib}
\IfFileExists{neurips_2026.sty}{\usepackage[preprint]{neurips_2026}}{}

\usepackage[utf8]{inputenc}
\usepackage[T1]{fontenc}
\usepackage{url}
\usepackage{booktabs}
\usepackage{amsfonts}
\usepackage{nicefrac}
\usepackage{microtype}
\usepackage{amsmath}
\usepackage{amssymb}
\usepackage{mathtools}
\usepackage{amsthm}
\usepackage[table]{xcolor}  % [table] loads colortbl -> \arrayrulecolor (grid-style tables)
\usepackage{array}
\usepackage{graphicx}
\usepackage{float}
\usepackage{enumitem}
\usepackage{longtable}
\usepackage{multirow}
\usepackage[hidelinks,colorlinks=true,citecolor=blue,linkcolor=blue,urlcolor=blue]{hyperref}
\usepackage[capitalize,noabbrev]{cleveref}
\usepackage{subcaption}
\usepackage{pifont}
\newcommand{\cmark}{\ding{51}}
\newcommand{\xmark}{\ding{55}}
% Provenance-name marker (house convention): every reference to a named experiment of this
% campaign (\runref{exp 010}, \runref{exp 020}, ...) is set bold inside square brackets.
\newcommand{\runref}[1]{\textbf{[#1]}}

% Grid-style hyperparameter tables (shared skill rnd-experiment-tex-track).
\setlength{\arrayrulewidth}{0.5pt}                 % \hline/\vline only; booktabs rules unaffected
\definecolor{tabgrey}{gray}{0.65}
\newcommand{\greyvline}{{\color{tabgrey}\vline}}   % braces keep the color inside the rule
\newcolumntype{?}{!{\greyvline}}                   % column separator: one grey vertical rule
\newcommand{\greyrule}{\arrayrulecolor{tabgrey}\hline\arrayrulecolor{black}} % within a category
\newcommand{\blackrule}{\arrayrulecolor{black}\hline}                        % category boundary

\mathtoolsset{showonlyrefs}
\DeclareMathOperator*{\argmax}{arg\,max}
\DeclareMathOperator*{\argmin}{arg\,min}

\theoremstyle{plain}
\newtheorem{theorem}{Theorem}[section]
\newtheorem{proposition}[theorem]{Proposition}
\theoremstyle{definition}
\newtheorem{definition}[theorem]{Definition}
\theoremstyle{remark}
\newtheorem{remark}[theorem]{Remark}

\title{Uniform $1/\sqrt{n}$ Bonus Decay at Every Position:\\ Development Document}

\author{%
  Anonymous Author(s)\\
  Affiliation\\
  \texttt{email}
}

\begin{document}

\maketitle

\tableofcontents
\clearpage

%==============================================================================
\section{Goal}
\label{sec:goal}

Random Network Distillation reads its exploration bonus off the residual between a trainable
predictor and a frozen random target. The count-based reading of that bonus treats it as a
stand-in for $\min(1, 1/\sqrt{n})$ at a state visited $n$ times, and the convergence-rate runs
of the parent project measured how far the real bonus is from that reading: with the best
learning-rate schedule the \emph{aggregate} decay matched $n^{-1/2}$ almost exactly
(fitted slope $-0.503$), but different maze positions started from different values and
decayed at visibly different rates. This campaign's goal is a training method whose
per-position bonus $b_i(n)$ satisfies, at \emph{every} position $i$ simultaneously:
\begin{enumerate}[leftmargin=*]
\item \textbf{Start at one.} $b_i(0) = 1$ at every position --- exactly, or through a
  normalization that an online implementation could compute (no oracle knowledge, no
  knowledge of the evaluation set).
\item \textbf{Decay as one over the square root of the visit count.}
  $b_i(n) \approx \min\!\bigl(1, m_i(n)^{-1/2}\bigr)$, where $m_i(n)$ is the number of times
  position $i$ has been trained on after $n$ update steps. In the uniform full-batch regime
  $m_i(n) = n$ for every $i$; under non-uniform visitation each position must follow its
  \emph{own} count.
\end{enumerate}
The deliverable is an algorithm — preferably a neural network — that plugs into a Random
Network Distillation pipeline and delivers this decay law, with generalization beyond the
maze demonstrated on AntMaze observations and Atari frames. Generalization is evaluated the
same way as everything else in this document: by the convergence rate of the bonus against
the count oracle's curve on fixed observation sets --- never by end-to-end reinforcement
learning, which is outside this campaign's scope.

The campaign ran as an autonomous experiment loop in the style of the
\texttt{autoresearch} repository\footnote{\url{https://github.com/karpathy/autoresearch},
cloned into the project for reference.}: one editable method file, a fixed measurement
protocol and metric, one commit and one ledger row per experiment. Fifty-plus experiments
later, the answer has four layers: a precise account of \emph{why} ordinary optimizers
cannot deliver the law (\Cref{sec:mechanism}); an exact construction that solves the
uniform regime to the achievable limit and a statistical construction that tracks
per-position counts under arbitrary visitation
(\Cref{sec:iteration-experiments,sec:nonuniform,sec:capacity}); and the neural forms of
both, evaluated on the high-dimensional domains (\Cref{sec:neural}).

%==============================================================================
\section{Notation}
\label{sec:notation}

\begin{table}[H]
\centering
\small
\renewcommand{\arraystretch}{1.25}
\begin{tabular}{@{}l l@{}}
\toprule
\textbf{Symbol} & \textbf{Meaning} \\
\midrule
$x_i$ & evaluation position $i$ (input $[x, y, 0, 0]$; \Cref{sec:platform} lists the sets) \\
$P$ & number of positions in a point set \\
$n$ & update-step index ($0$ = at initialization) \\
$m_i(n)$ & number of times position $i$ was in a training batch during steps $1 \ldots n$ \\
$f$ & frozen random target network \\
$g_n$ & trained predictor after $n$ steps ($g_0$ = at initialization) \\
$b_i(n)$ & the method's bonus readout at position $i$ after $n$ steps \\
$T_i(n)$ & target curve $\min\!\bigl(1, m_i(n)^{-1/2}\bigr)$ (the count oracle's bonus) \\
$K$ & tangent-kernel Gram matrix over a point set; eigenvalues $\lambda_j$, eigenvectors $v_j$ \\
$\eta_t$ & learning rate (or shrink factor) at step $t$ \\
$\phi(x)$ & feature vector of a linear-head method; $\hat\phi = \phi / \lVert \phi \rVert_2$ \\
$d$ & output dimension of the predictor (or of the coin vector) \\
$S$ & number of seeds of an experiment (each seed re-draws every initialization) \\
\bottomrule
\end{tabular}
\caption{Symbols used throughout. Machine timestamps in run records are Eastern; every time
displayed in this document is Pacific, marked PT.}
\label{tab:notation}
\end{table}

%==============================================================================
\section{Metrics}
\label{sec:metrics}

Define the fit window as the checkpoint steps $2 \le n \le 4096$ (the parent project's
burn-in convention: step $0$ is the starting value, step $1$ is excluded, nothing else is
trimmed). Define, for seed $s$ and position $i$, the per-position deviation
\begin{align}
\label{eq:dev}
\mathrm{dev}_{s,i}
\;=\;
\sqrt{\;\frac{1}{|\mathcal{W}|}\sum_{n \in \mathcal{W}}
\bigl(\log b_{s,i}(n) - \log T_i(n)\bigr)^2\;},
\end{align}
where $\mathcal{W}$ is the set of in-window checkpoint steps and $T_i(n) = \min(1,
m_i(n)^{-1/2})$. The deviation is a root mean square in log space, so it penalizes a wrong
level and a wrong slope together, and a method scores well only if its curve sits \emph{on}
the count oracle's curve, not merely parallel to it. The campaign's metrics, in ranking
order:
\begin{enumerate}[leftmargin=*]
\item \textbf{dev\_worst} $= \frac{1}{S}\sum_s \max_i \mathrm{dev}_{s,i}$ --- the primary
  metric. ``Every position at the same time'' means the worst position is the score.
\item \textbf{dev\_mean} $= \operatorname{mean}_{s,i} \mathrm{dev}_{s,i}$ --- the tiebreak.
\item \textbf{start\_dev} $= \sqrt{\operatorname{mean}_{s,i} (\log b_{s,i}(0))^2}$ ---
  requirement (1) in one number ($0$ exactly when every curve starts at $1$).
\item \textbf{The slope battery.} Every per-position curve (seed-mean) is fitted with the
  parent project's floored power model $y(n) = c + a\,(n + n_0)^{-\alpha}$ (nonlinear least
  squares on $\log y$, bounded, multi-start); the reported slope is $-\alpha$. The battery
  reports the mean, standard deviation, minimum and maximum of the per-position slopes, and
  the aggregate-curve slope in both the raw and the normalize-then-average convention of the
  parent project (so the numbers are comparable to its $-0.503$).
\end{enumerate}
Two floors bound what any method can score. A method whose state cannot distinguish ``never
visited'' from ``visited once'' produces $(m+1)^{-1/2}$ instead of $\min(1, m^{-1/2})$; over
this window that off-by-one contributes a deviation of $0.0381$, identical at every position.
A method that estimates the count statistically from $d$ independent draws carries a
chi-type fluctuation of relative standard deviation about $\sqrt{1/(2d)}$ on the readout
norm ($\approx 3\%$ at $d = 512$), which enters dev\_worst through the maximum over
positions.

%==============================================================================
\section{Why the decay rate depends on the position}
\label{sec:mechanism}

\subsection{Mechanism one: optimization against a fixed target}
\label{sec:mechanism-one}

Recall $K$ is the tangent-kernel Gram matrix over the point set, with eigenvalues
$\lambda_j$ and eigenvectors $v_j$, and $\eta_t$ the learning rate. In the linearized
regime, full-batch gradient descent on the squared distillation loss evolves the residual
vector $r$ as $r_{n} = \prod_{t \le n}(I - \eta_t K)\, r_{0}$. Define the mixture weight
$w_{ij} = v_{j,i}^2$ (the squared $i$-th entry of eigenvector $j$); each row of this
evolution gives the per-position law
\begin{align}
\label{eq:per-position-law}
b_i(n)^2 \;\approx\; \sum_j w_{ij}\, r_j(0)^2 \prod_{t \le n} (1 - \eta_t \lambda_j)^{2},
\qquad \sum_j w_{ij} = 1 ,
\end{align}
so each position carries its own probability distribution over the eigenvalues while the
aggregate uses the uniform mixture --- which is precisely how an aggregate slope can look
clean while positions disagree. Three consequences:
\begin{enumerate}[label=(\roman*),leftmargin=*]
\item \textbf{A constant rate gives geometric decay, then saturation.} Each mode contracts
  by $(1 - \eta\lambda_j)$ per step; fast modes vanish, the residual settles onto the
  flattest surviving mode and the curve flattens. This is the parent project's constant-rate
  Adam observation (aggregate slope about $-0.8$, then saturating).
\item \textbf{A $1/t$ schedule gives every mode its own power law.} With $\eta_t = a/t$,
  mode $j$ decays as $n^{-a\lambda_j}$: the exponent is \emph{proportional to the
  eigenvalue}. A position's curve is a mixture of power laws, its fitted slope depends on its
  eigenvector mixture and drifts with the window. The parent project's aggregate $-0.503$ is
  the schedule constant placing the spectrum's bulk near $a\bar\lambda \approx 1/2$ --- a
  calibration, not a mechanism.
\item \textbf{No global schedule can equalize the positions.} For any schedule, mode $j$'s
  log-decay is $\approx -\lambda_j \sum_{t\le n}\eta_t$: the \emph{ratio} of two modes'
  log-decays is pinned at $\lambda_j / \lambda_k$ whatever $\{\eta_t\}$ is. Equal
  per-position rates therefore require either a flat spectrum ($K \propto I$), delocalized
  eigenvectors ($w_{ij}$ identical for all $i$), or a step that does not pass through $K$ at
  all. The measured tangent kernel of this project's $4 \to 256 \to \mathrm{ReLU} \to 128$
  network on a $10\times10$ grid has condition number $5.6 \times 10^{5}$, so the spectrum
  is very far from flat.
\end{enumerate}
The linearized picture is checked directly in \Cref{fig:linearization-check}: simulating the
linearized dynamics of the real network (matrix-free Jacobian products, same schedule, same
checkpoints, same floored fit) predicts the measured per-position slopes of \runref{exp 002}
with correlation $0.75$ and a matching range --- the residual heterogeneity is the tangent
kernel, not noise.

\begin{figure}[H]
\centering
\includegraphics[width=0.5\linewidth]{figures/linearization_check.pdf}
\caption{Per-position slopes predicted by the linearized dynamics (one seed, matrix-free
$J J^\top$ products at initialization) against the measured slopes of \runref{exp 002}
(seed-mean over $10$ seeds), cell midpoints, SGD-$1/t$. The prediction uses a single seed
while the measurement averages ten, and real training leaves the linearized regime late in
the run, which bounds the agreement.}
\label{fig:linearization-check}
\end{figure}

Two further facts locate the heterogeneity in space. A bias-free ReLU network is positively
homogeneous, so the kernel diagonal grows with the input radius ($K_{ii} \propto
\lVert x_i \rVert^2$ up to a constant): the maze corners hold the largest inputs, train
fastest, and were exactly the worst positions in the AdaGrad diagnostics of
\Cref{sec:iteration-experiments}. And the same homogeneity makes the \emph{initial} bonus
position-dependent --- in the wide limit $\mathbb{E}\lVert g_0(x) - f(x)\rVert^2$ is
proportional to the kernel diagonal --- which is requirement (1)'s obstacle: the untrained
bonus field is a bowl centered near the coordinate origin.

\subsection{Mechanism two: statistical averaging of fresh random targets}
\label{sec:mechanism-two}

If the target presented at position $i$ on each visit is a \emph{fresh} independent draw
with mean $\mu(x_i)$ and per-coordinate variance $\sigma^2$ (rather than a fixed value), the
least-squares optimum at a position visited $m_i$ times is that position's own running mean
of draws, and the expected squared residual is $\sigma^2 d / m_i$ \emph{exactly} --- an
identity about sample means, involving no curvature, no learning rate, and no other
position. The residual norm then decays as $m_i^{-1/2}$ with the same exponent and the same
constant at every position, using each position's own count, under any visitation pattern.
This is the coin-flip pseudo-count construction of Lobel, Bagaria and Konidaris (ICML 2023);
\Cref{sec:iteration-experiments} builds its exact-least-squares version. The cost is a
statistical fluctuation floor (the chi-type noise of \Cref{sec:metrics}) and a
representation requirement: the function class must be able to hold a distinct running mean
per position, which \Cref{sec:iteration-experiments} shows is a real constraint, not a
formality.

\subsection{The three knobs}
\label{sec:three-knobs}

Every method the campaign tested decomposes into three independent choices:
\begin{enumerate}[leftmargin=*]
\item \textbf{Who shares a rate} --- whether the update passes through the tangent kernel
  (per-mode rates, position-dependent), through a function-space step (one shared rate), or
  through per-position statistics (each position its own count).
\item \textbf{The shape} --- what turns the shared rate into $n^{-1/2}$: a shrink schedule
  whose cumulative product is $n^{-1/2}$; a per-visit map applied to the bonus value itself;
  the sample-mean identity; or a quartic loss whose constant-step dynamics integrate to
  $n^{-1/2}$.
\item \textbf{The level} --- what makes the curve start at $1$: a frozen copy of the initial
  predictor as a per-position denominator; a zero-initialized head with a unit-normalized
  target or prior; or a Rademacher draw whose norm is $\sqrt{d}$ exactly.
\end{enumerate}

%==============================================================================
\section{Experiment platform}
\label{sec:platform}

The campaign adapts the \texttt{autoresearch} loop: a fixed harness owns the protocol and
the metrics; one editable method file owns everything about how the bonus is produced
(networks, initialization, optimizer, loss, readout, auxiliary state); each experiment is
one edit, one commit, one Slurm job, one ledger row. The loop may not modify the harness,
may not read the step counter or visit counts in the \emph{readout} (a learning-rate
schedule may use the step counter --- the optimizer may know the time, the readout may not),
and may not key a lookup table by training-set index (methods must be functions of the input
coordinates). One experiment runs its point-set--seed cells as parallel single-thread
workers and finishes in about half a minute, which is what made a thirty-experiment campaign
practical in one night.

\begin{table}[H]
\centering
\footnotesize
\setlength{\tabcolsep}{3.5pt}
\renewcommand{\arraystretch}{1.2}
\begin{tabular}{@{}>{\raggedright\arraybackslash}p{3.6cm} ?>{\raggedright\arraybackslash}p{9.6cm}@{}}
\blackrule
\textbf{Protocol knob} & \textbf{Value (fixed for every experiment)} \\
\blackrule
\multicolumn{2}{@{}l}{\textit{Point sets (evaluation = training support)}} \\
\greyrule
\texttt{cell\_\allowbreak midpoints} & the midpoint of every cell of the $9 \times 12$ maze, walls included; $108$ points (primary set) \\
\greyrule
\texttt{center\_\allowbreak square} & $10\times10$ sub-cell midpoints of the $1\times1$ square at the maze center, spacing $0.1$; $100$ points \\
\greyrule
\texttt{top\_\allowbreak right\_\allowbreak cell} & the same discretization of the top-right goal cell; $100$ points (validation runs only) \\
\greyrule
\texttt{dense\_\allowbreak grid} & $2\times2$ sub-cell midpoints of every cell; $432$ points --- more positions than a $256$-wide feature layer (capacity stress, added mid-campaign) \\
\greyrule
\texttt{antmaze\_\allowbreak states} & $246$ real AntMaze\_Large observation vectors ($29$ dimensions: $x$, $y$, plus the $27$ contact-force-free core dimensions), collected by a seeded random policy and stratified over $41$ maze cells (generalization domain, phase 2) \\
\greyrule
\texttt{atari\_\allowbreak frames} & $512$ MontezumaRevenge frames from seeded random play, preprocessed as RND's predictor sees them (grayscale, $84\times84$, single frame) --- the image-input generalization domain (phase 2) \\
\blackrule
\multicolumn{2}{@{}l}{\textit{Training regimes}} \\
\greyrule
\texttt{uniform\_\allowbreak fullbatch} & one update per step on the full point set; $m_i(n) = n$ \\
\greyrule
\texttt{nonuniform} & each step trains on $32$ points sampled i.i.d.\ from a fixed law spanning one decade of probability across the maze's $x$ extent; targets use realized counts \\
\greyrule
\texttt{heldout\_\allowbreak uniform} & full-batch training on a fixed $80\%$ of the point set; the held-out $20\%$ are never trained, their counts stay $0$, and the metric scores whether their bonus stays at the oracle value $1$ --- generalization measured inside the convergence framework (phase 2) \\
\blackrule
\multicolumn{2}{@{}l}{\textit{Measurement}} \\
\greyrule
steps, checkpoints & $4096$ steps; readout stored at step $0$ and the $79$ log-spaced steps $\mathrm{round}(2^{j/8})$ \\
\greyrule
input & raw $[x, y, 0, 0]$, no observation normalization \\
\greyrule
seeds & $10$ per iteration experiment, $30$ per validation run \\
\greyrule
metrics & \Cref{sec:metrics}; records atomically checkpointed, diverged cells excluded and counted \\
\blackrule
\end{tabular}
\caption{The fixed protocol. Grid-style table: grey rules separate rows within a category,
black rules separate categories.}
\label{tab:protocol}
\end{table}

\clearpage
%==============================================================================
\section{Iteration experiments}
\label{sec:iteration-experiments}

Thirty-three experiments ran under the fixed protocol on 2026-08-17/18 (PT); every experiment
is one row of \Cref{tab:ledger} and one commit of the method file. This section walks the
five method families in the order the campaign discovered them; \Cref{sec:nonuniform} and
\Cref{sec:capacity} then stress the winners. Experiment numbers in brackets
(\runref{exp 010}) name the ledger rows.

\subsection{Baselines: the optimizer family}
\label{sec:baselines}

Four experiments established the board. Constant-rate Adam \runref{exp 001} reproduces the
parent project's finding at $10$ seeds: slope mean $-0.86$, too fast everywhere, and an
initial bonus spread of $1.37$ in log (the untrained bowl). The parent project's best
schedule, SGD with $\eta_t = \eta_0/(1 + t/t_0)$ \runref{exp 002}, exposes what its
aggregate $-0.503$ hid: the normalized aggregate slope is $-0.80$ on the cell midpoints but
$-0.38$ on the center square --- the pooled aggregate averaged heterogeneous point sets.
Normalizing the readout by a frozen copy of the initial predictor \runref{exp 003} sets
start\_dev to exactly $0$ (requirement (1) is that cheap: one stored network) while leaving
every slope unchanged, cleanly separating the level problem from the rate problem.

AdaGrad is the best of the family, consistent with its cumulative (never-forgetting)
accumulator: at learning rate $10^{-2}$ \runref{exp 007} the aggregate slopes are $-0.51$
and $-0.50$ on the two point sets with per-position spread $0.11$; at $3\times10^{-3}$
\runref{exp 009} dev\_mean drops to $0.32$, the best of any pure-gradient method. Its two
defects are visible in \Cref{fig:adagrad-diagnostics}: every curve falls to about $0.55$
on the very first step (the oracle stays at $1$ until the first visit), and the maze
corners --- the largest-radius inputs, with the largest kernel diagonal
(\Cref{sec:mechanism-one}) --- decay about $0.17$ too fast.

\begin{figure}[H]
\centering
\includegraphics[width=\linewidth]{figures/adagrad_diagnostics.pdf}
\caption{AdaGrad diagnostics on the cell midpoints (\runref{exp 007}, seed-mean). Left to
right: per-position deviation over the maze (hatched cells are walls), per-position fitted
slope, the six worst positions' curves against $\min(1, n^{-1/2})$ (grey), and the
per-position slope histogram. The worst positions are the four corners; the first-step drop
and the too-steep tail are the two visible defects.}
\label{fig:adagrad-diagnostics}
\end{figure}

\subsection{Global schedules and averaging cannot fix the spread}
\label{sec:global-attempts}

Four experiments tested repairs that keep the update passing through the tangent kernel;
all four fail in the way \Cref{sec:mechanism-one} predicts.
\begin{enumerate}[label=(\roman*),leftmargin=*]
\item \textbf{Polyak-averaged readout} \runref{exp 004}: constant-step SGD with a uniformly
  averaged parameter copy and a square-root readout. Averaging does make every mode's
  asymptotic decay $1/n$, but each mode enters that regime only after its own transient
  $n \gtrsim 1/(\eta\lambda_j)$; within the window the center square barely decays
  (aggregate $-0.10$) while the cell midpoints reach $-0.46$.
\item \textbf{Inverse-squared-error loss reweighting} \runref{exp 006}: weighting each
  sample's loss by the inverse of its own detached squared error should equalize relative
  progress, but the weights act through the kernel's off-diagonals and the spread survives
  (slope standard deviation $0.69$).
\item \textbf{Per-sample shrink targets fitted by inner gradient descent} \runref{exp 013}:
  the mean slope lands on $-0.51$ but the slow kernel modes stay under-shrunk (center-square
  aggregate $-0.22$) --- inner gradient fitting cannot realize a functional step on modes it
  cannot reach, which is the same spectrum problem one level down.
\item \textbf{Quartic loss} \runref{exp 030}: minimizing the fourth power of the normalized
  residual norm at a constant step has the closed-form diagonal dynamics
  $b(n) = (2\eta\kappa n + b(0)^{-2})^{-1/2}$ --- exponent $-1/2$ by construction, no
  schedule, no counter, and the initial value is forgotten. The coupled kernel slows the
  real network to a slope nearer $-0.4$ with a two-fold level spread; it is the best
  \emph{schedule-free pure-gradient} approximation, and not more.
\end{enumerate}
One failed experiment earned its place in the mechanism story: realizing a functional
shrink with inner \emph{Adam} steps \runref{exp 005} produced slope $-1.8$ --- Adam's
step magnitude does not scale with the distance to the (very close) inner target, so each
inner loop overshoots it toward the frozen target and the method degenerates to fast
distillation. Inner fitting must contract (gradient steps proportional to the residual), or
be replaced by a solve.

\subsection{The exact construction: a residual-encoded per-visit map, realized by a solve}
\label{sec:exact-construction}

The count oracle's curve satisfies a one-step recurrence that never mentions the count.
Define $r$ as the current normalized bonus at a position; if $r = (m+1)^{-1/2}$ after $m$
visits, then after one more visit
\begin{align}
\label{eq:shrink-map}
r' \;=\; \frac{r}{\sqrt{1 + r^2}} \;=\; (m+2)^{-1/2},
\end{align}
so the map $r \mapsto r/\sqrt{1+r^2}$, applied once per visit, walks $1 \to 1/\sqrt{2} \to
1/\sqrt{3} \to \cdots$ \emph{The bonus value itself is a sufficient statistic for the visit
count}: no counter, no clock, valid under any visitation pattern.

\runref{exp 010} realizes the map exactly. The predictor is a linear head $W$ on the frozen
first layer's features $\phi(x)$; one update computes each visited sample's normalized
residual $r_i$, forms the shrunk residual targets $s_i e_i$ with $s_i = 1/\sqrt{1+r_i^2}$,
and applies the minimum-Frobenius-change head update that interpolates the batch constraints
exactly:
\begin{align}
\label{eq:min-change}
\Delta W = E^\top G^{-1} \Phi,
\qquad
G = \Phi\Phi^\top + 10^{-8} I,
\qquad
E_i = (s_i - 1)\, e_i ,
\end{align}
with $\Phi$ the batch feature matrix and $e_i$ the residual vectors. On the uniform
benchmark the result is the achievable limit: fitted slope $-0.500000$ at every one of the
$208$ positions (standard deviation $4\times10^{-6}$ over positions), start exactly $1$, and
deviation $0.0381$ = the off-by-one floor of \Cref{sec:metrics}, identical at every
position. In the three-knob language of \Cref{sec:three-knobs}: the solve decides who shares
a rate (everyone), the map decides the shape, the initial-copy denominator decides the
level.

Two follow-ups located the construction's strength and its weakness. On the dense grid ---
$432$ positions against $256$ features, where exact interpolation is impossible ---
\runref{exp 023} still scores $0.0383$ with slope spread $4\times10^{-5}$: the map reads the
\emph{achieved} state, so whatever the solve under-delivers on one visit is corrected on the
next; the construction is self-correcting under capacity stress. Under non-uniform
minibatch visitation, however, it fails completely (\Cref{sec:nonuniform}).

\subsection{The statistical construction: coin flips with an exact least-squares head}
\label{sec:statistical-construction}

The coin-flip construction (mechanism two) replaces the fixed target by a fresh Rademacher
vector $c \in \{-1,+1\}^d$ per visit and reads the bonus off the norm of the learned
running mean. Three campaign choices make it exact at the edges: the head is solved by
ridge least squares on accumulated normal equations (so the optimization gap of the
gradient-descent version is zero); the head starts at zero with a frozen prior network
unit-normalized per input ($\lVert p(x)\rVert = \sqrt{d}$ exactly), so the bonus is exactly
$1$ at every never-visited input; and a Rademacher first draw has norm $\sqrt{d}$ exactly,
so the bonus is also exactly $1$ after the first visit --- matching $\min(1, m^{-1/2})$ at
$m = 0$ and $1$ with no off-by-one.

The first attempt \runref{exp 014} failed for a reason worth recording: the project's
default zero-bias ReLU features are homogeneous half-plane hinges whose span on the $108$
cell midpoints has numerical rank $106$ and condition number $10^{10}$ --- the least-squares
optimum cannot reach per-position running means, and the readout saturates on the
unrepresentable component. Swapping in $256$ Gaussian radial-basis features (bandwidth
$0.75$, Gram condition about $10^{2}$) \runref{exp 016} restored the construction on the
cell midpoints (aggregate slope $-0.5013$) but under-resolved the $0.1$-spaced center
square. The fix that closed both gaps is a \emph{growing, data-dependent dictionary}
\runref{exp 019}: every visited state farther than $0.05$ from the dictionary becomes a new
center whose bandwidth is half its nearest-center distance at insertion, and the
normal-equation state grows with it. Resolution concentrates where states are visited, at
the spacing they are visited at; the uniform benchmark lands at dev\_worst $0.053$ ---
essentially the chi fluctuation floor at $d = 512$.

\subsection{The closed-form count: an elliptical posterior readout}
\label{sec:elliptical-readout}

The deterministic ideal of the coin-flip family drops the coins entirely: keep the running
second-moment matrix $A = I + \sum_{\text{visits}} \hat\phi\hat\phi^\top$ of unit-normalized
features on the same adaptive dictionary and read out the posterior standard deviation
\begin{align}
\label{eq:elliptical}
b(x) \;=\; \sqrt{\hat\phi(x)^\top A^{-1} \hat\phi(x)} ,
\end{align}
which is exactly $1$ before any visit and exactly $(1+m_i)^{-1/2}$ for mutually orthogonal
features --- the elliptical bonus of the parent project, repurposed as the bonus itself.
\runref{exp 027} measures dev\_worst $0.202$ (uniform): the slope is near-perfect
($-0.499 \pm 0.002$) and the entire remaining deviation is neighboring-cell feature overlap
inflating effective counts --- generalization itself, showing up as a level shift. The
bandwidth cap is the dial between per-position exactness and spatial generalization
(\runref{exp 032}/\runref{exp 033} probe a tighter cap).

% >>> AUTO-GENERATED TABLE START: ledger
{\footnotesize
\begin{longtable}{@{}r >{\raggedright\arraybackslash}p{5.2cm} r r r r r@{}}
\caption{The campaign ledger: every iteration experiment, in submission order. Columns are the primary and secondary deviation metrics, the start deviation, the per-position slope mean and standard deviation (\Cref{sec:metrics}); (discarded) marks abandoned directions. Numbers are the experiment's own 10-seed metrics; the exact method per row is the method file at the ledger's commit.}\label{tab:ledger}\\
\toprule
exp & method / finding & dev\_worst & dev\_mean & start & slope & sl.\ std \\
\midrule\endfirsthead
\toprule
exp & method / finding & dev\_worst & dev\_mean & start & slope & sl.\ std \\
\midrule\endhead
001 & adam baseline (constant Adam 1e-3, raw l2 readout) & 1.618 & 1.214 & 1.37 & -0.862 & 0.133 \\
002 & sgd1t baseline (eta0 1e-2, t0 1e2) & 2.471 & 1.918 & 1.37 & -0.969 & 0.776 \\
003 & sgd1t + initial-copy readout norm & 2.634 & 1.407 & 0.00 & -0.867 & 0.564 \\
004 & const SGD + Polyak-averaged sqrt readout & 2.541 & 2.025 & 0.00 & -0.453 & 0.441 \\
005 & functional shrink w/ sqrt schedule (discarded) & 1.686 & 0.776 & 0.00 & -1.801 & 0.508 \\
006 & SGD-1/t + inverse-squared-error reweighting (discarded) & 2.393 & 1.802 & 0.00 & -0.765 & 0.689 \\
007 & AdaGrad + initial-copy readout & 1.227 & 0.761 & 0.00 & -0.568 & 0.106 \\
008 & sphere-projected inputs + AdaGrad (discarded) & 1.585 & 1.188 & 0.00 & -1.043 & 0.598 \\
009 & AdaGrad lr 3e-3 & 1.609 & 0.323 & 0.00 & -0.736 & 0.201 \\
010 & linear head + exact min-change interpolation + residual-en\,\ldots & 0.038 & 0.038 & 0.00 & -0.500 & 0.000 \\
011 & = exp-010 method under NONUNIFORM visitation & 1.467 & 0.809 & 0.00 & -0.500 & 0.000 \\
012 & AdaGrad 3e-3 under NONUNIFORM & 1.412 & 0.676 & 0.00 & -0.782 & 0.217 \\
013 & MLP + per-sample shrink targets + inner SGD (discarded) & 2.402 & 1.170 & 0.00 & -0.514 & 0.283 \\
014 & coin-flip RLS on zero-bias ReLU features, uniform (discarded) & 2.535 & 1.247 & 0.00 & -0.901 & 0.267 \\
015 & same under nonuniform (discarded) & 1.905 & 0.780 & 0.00 & -1.201 & 0.994 \\
016 & coin-flip RLS on RBF features (sigma 0.75), uniform & 1.008 & 0.333 & 0.00 & -0.574 & 0.126 \\
017 & same under NONUNIFORM & 0.885 & 0.350 & 0.00 & -0.635 & 0.109 \\
018 & shrink on fixed RBF, NONUNIFORM & 1.076 & 0.588 & 0.00 & -0.512 & 0.074 \\
019 & coin-flip + adaptive dictionary, uniform & 0.053 & 0.030 & 0.00 & -0.503 & 0.004 \\
020 & same, NONUNIFORM & 0.455 & 0.143 & 0.00 & -0.624 & 0.106 \\
021 & same, dense\_grid 432 points & 0.059 & 0.030 & 0.00 & -0.503 & 0.004 \\
022 & coin-flip fixed-RBF on dense\_grid & 1.397 & 0.271 & 0.00 & -0.555 & 0.055 \\
023 & shrink + global ReLU on dense\_grid (432 > 256 features) & 0.038 & 0.038 & 0.00 & -0.500 & 0.000 \\
024 & shrink + adaptive dictionary, uniform & 0.038 & 0.038 & 0.00 & -0.500 & 0.000 \\
025 & same, NONUNIFORM & 1.310 & 0.522 & 0.00 & -0.576 & 0.199 \\
026 & same, dense\_grid & 0.038 & 0.038 & 0.00 & -0.500 & 0.000 \\
027 & closed-form elliptical posterior on adaptive dictionary, u\,\ldots & 0.202 & 0.054 & 0.00 & -0.499 & 0.002 \\
028 & same, NONUNIFORM & 0.535 & 0.258 & 0.00 & -0.564 & 0.080 \\
029 & elliptical adaptive on dense\_grid & 0.126 & 0.048 & 0.00 & -0.499 & 0.001 \\
030 & quartic loss constant step, uniform & 1.015 & 0.692 & 0.00 & -0.398 & 0.047 \\
031 & same, NONUNIFORM & 0.837 & 0.506 & 0.00 & -0.349 & 0.043 \\
032 & elliptical with bandwidth cap 0.35, uniform & 0.114 & 0.044 & 0.00 & -0.500 & 0.001 \\
033 & same, NONUNIFORM & 0.478 & 0.228 & 0.00 & -0.576 & 0.084 \\
034 & scrambled-Hadamard coin sequences on the adaptive dictionary & 0.000 & 0.000 & 0.00 & 0.000 & 0.000 \\
035 & Hadamard coins + per-block column permutation (val\_114-115) & 0.314 & 0.008 & 0.00 & -0.501 & 0.001 \\
036 & collision-free block spikes & 0.033 & 0.003 & 0.00 & -0.500 & 0.000 \\
? & abl 201 coin-flip d=128 & 0.108 & 0.060 & 0.00 & -0.506 & 0.000 \\
? & abl 202 coin-flip insertion radius 0.2 (uniform) & 1.062 & 0.531 & 0.00 & -0.701 & 0.000 \\
? & abl 203 same, nonuniform & 0.994 & 0.435 & 0.00 & -0.761 & 0.000 \\
? & abl 204 elliptical bandwidth cap 0.2 & 0.077 & 0.043 & 0.00 & -0.500 & 0.000 \\
? & hor 301 shrink at 32768 steps & 0.033 & 0.033 & 0.00 & -0.500 & 0.000 \\
? & hor 303 AdaGrad 3e-3 at 32768 steps & 1.296 & 0.322 & 0.00 & -0.631 & 0.000 \\
? & hor 302 coin-flip adaptive at 32768 steps & 0.049 & 0.030 & 0.00 & -0.502 & 0.003 \\
? & val 118 collision-free Hadamard on dense\_grid, 30 seeds & 0.054 & 0.003 & 0.00 & -0.500 & 0.000 \\
037 & gradient CFN, Adam 1e-4 (paper rate), maze+antmaze (discarded) & 2.202 & 1.849 & 0.00 & -0.648 & 0.347 \\
038 & gradient CFN Adam 1e-3 (discarded) & 1.424 & 1.123 & 0.00 & -0.790 & 0.230 \\
039 & gradient CFN Adam 1e-2 (discarded) & 1.148 & 1.159 & 0.00 & -0.696 & 0.264 \\
040 & gradient CFN AdaGrad 1e-2 (discarded) & 1.469 & 1.037 & 0.00 & -0.395 & 0.088 \\
041 & gradient CFN AdaGrad 1e-1 (discarded) & 1.766 & 1.654 & 0.00 & -1.402 & 0.763 \\
042 & gradient CFN SGD-1/t eta0 1e-2 (discarded) & 2.057 & 1.541 & 0.00 & -0.598 & 0.146 \\
043 & gradient CFN SGD-1/t eta0 1e-1 (discarded) & 1.856 & 1.455 & 0.00 & -1.065 & 0.427 \\
048 & deep random features + exact last-layer shrink head, maze+\,\ldots & 0.038 & 0.038 & 0.00 & -0.500 & 0.000 \\
049 & deep features + exact coin-flip head, maze+antmaze pooled & 0.054 & 0.030 & 0.00 & -0.503 & 0.004 \\
\bottomrule
\end{longtable}}
% >>> AUTO-GENERATED TABLE END: ledger

\begin{figure}[H]
\centering
\includegraphics[width=\linewidth]{figures/champion_curves.pdf}
\caption{Per-position bonus curves on the cell midpoints (uniform full batch, seed-mean; one
thin line per position, $108$ per panel; grey line $= \min(1, n^{-1/2})$). Left to right:
constant-rate Adam \runref{exp 001} (raw readout: spread starts and too-fast decay), AdaGrad
$3\times10^{-3}$ with the initial-copy readout \runref{exp 009} (starts at $1$, first-step
drop and corner spread remain), the residual-encoded shrink \runref{exp 010} (every curve on
the oracle), and coin-flip counting on the adaptive dictionary \runref{exp 019} (every curve
on the oracle up to the chi fluctuation).}
\label{fig:champion-curves}
\end{figure}

%==============================================================================
\section{Non-uniform visitation}
\label{sec:nonuniform}

The nonuniform regime samples $32$ positions per step from a law spanning one decade of
probability across the maze, and scores each position against its \emph{own} realized count.
This is where the two exact uniform-regime constructions separate:
\begin{enumerate}[label=(\roman*),leftmargin=*]
\item \textbf{The shrink map does not track counts} \runref{exp 011}. Every position's curve
  is the same aggregate curve (per-position slopes $-0.5002 \pm 0.0004$ against the
  \emph{global} step): the min-change solve moves the head globally, unvisited positions'
  residuals are dragged along through shared features, and the recurrence --- whose state is
  the corrupted bonus value itself --- integrates the interference. Localizing the features
  (\runref{exp 018}, \runref{exp 025}) reduces but does not remove the drift; the
  multiplicative recurrence amplifies whatever leaks through.
\item \textbf{Coin statistics track counts} \runref{exp 020}. The least-squares running mean
  is a batch statistic, not a recurrence: interference perturbs individual solves but
  averages out. With the adaptive dictionary the score is dev\_worst $0.455$, dev\_mean
  $0.143$ --- each position visibly following its own count.
\item \textbf{The closed-form elliptical readout is close behind} \runref{exp 028}
  (dev\_worst $0.535$): zero statistical noise, but the feature-overlap count inflation of
  \Cref{sec:elliptical-readout} costs more under heterogeneous counts.
\item \textbf{Gradient-descent baselines sit in between}: AdaGrad \runref{exp 012} is mildly
  count-sensitive (slope mean $-0.78$ against the global step) --- consistent with the
  density-dependent convergence times of the linearized theory --- and the quartic loss
  \runref{exp 031} inherits its uniform-regime level spread.
\end{enumerate}

\begin{figure}[H]
\centering
\includegraphics[width=0.9\linewidth]{figures/nonuniform_scatter.pdf}
\caption{Final bonus against final realized visit count, one dot per cell-midpoint position
(nonuniform regime, seed-mean, log--log; grey line $= m^{-1/2}$). The shrink construction
\runref{exp 011} collapses to a single bonus value regardless of the position's count; the
coin-flip construction \runref{exp 020} falls on the count oracle's line.}
\label{fig:nonuniform-scatter}
\end{figure}

%==============================================================================
\section{Capacity stress}
\label{sec:capacity}

The dense grid holds $432$ distinct positions against $256$-dimensional feature layers, so
exact interpolation across the whole set is impossible and every method must degrade
somehow. The measured degradations rank the mechanisms:
\begin{enumerate}[label=(\roman*),leftmargin=*]
\item The residual-encoded shrink \runref{exp 023} does not degrade at all
  (dev\_worst $0.0383$, slope spread $4\times10^{-5}$): the map's self-correction converts
  the solve's per-step realization error into a bounded wobble instead of a drift.
\item The adaptive dictionary \runref{exp 021}, \runref{exp 026} grows to $432$ centers and
  absorbs the stress ($0.059$ and $0.038$).
\item The fixed-bandwidth basis \runref{exp 022} hits the wall it cannot grow past
  (dev\_worst $1.40$): under-resolution is a resolution property, not a point-count
  property.
\end{enumerate}

\clearpage
%==============================================================================
\section{The neural algorithm and the generalization domains}
\label{sec:neural}

Phase 2 (a user redirection mid-campaign) sharpened the deliverable: the algorithm should be
a neural network that plugs into a Random Network Distillation pipeline, and its claim to
generality is tested on real high-dimensional observations --- the \texttt{antmaze\_states}
and \texttt{atari\_frames} domains and the \texttt{heldout\_uniform} regime of
\Cref{tab:protocol} --- always by convergence rate against the count oracle, never by
running reinforcement learning.

\subsection{Pure gradient training does not track the coin statistics}
\label{sec:neural-gradient}

The published coin-flip network trains an ordinary predictor toward fresh Rademacher targets
by gradient descent, and \Cref{sec:mechanism-two}'s statistics only pay off if training
tracks each state's running coin mean. Measured under this protocol --- one gradient step
per visit, exactly what an online bonus inside an RL loop would get --- it does not: across
seven optimizer variants (Adam at three rates including the paper's $10^{-4}$, AdaGrad at
two, SGD-$1/t$ at two; \runref{exp 037}--\runref{exp 043}) the pooled deviation never beats
$1.0$ and the per-state slopes scatter from $-0.4$ to $-1.4$. The mechanism of
\Cref{sec:mechanism-one} explains why: the update passes through the tangent kernel, so the
lag behind the moving least-squares optimum is spectrum-shaped and position-dependent. A
replay buffer (the paper's actual regime: stored (state, coin) pairs, four Adam minibatch
steps per visit, \runref{exp 052}) closes much of the lag and adds its own forgetting: the
ring's capacity bounds the effective count a state can accumulate.

\subsection{The RND-pluggable answer: a deep trunk with an exact last-layer head}
\label{sec:neural-lastlayer}

The construction that survives contact with every domain keeps the neural network for what
neural networks are good at --- representing observations --- and moves the counting into
the last layer, where it can be exact. A frozen deep trunk $\phi$ (the vector MLP or RND's
conv stack; random biases $N(0, 0.5^2)$ only for inputs of dimension $\le 8$, which need
them to break ReLU homogeneity) feeds one of the two phase-1 heads:
\begin{enumerate}[leftmargin=*]
\item \textbf{Shrink head} (\runref{exp 048}, \runref{exp 050}): bonus
  $r(x) = \lVert W\phi(x) - f(x)\rVert / \lVert f(x)\rVert$ with a zero-initialized head
  ($r \equiv 1$ before any visit) and the per-visit map $r \mapsto r/\sqrt{1+r^2}$ of
  \Cref{eq:shrink-map}, realized by the minimum-change solve of \Cref{eq:min-change}.
\item \textbf{Coin-flip head} (\runref{exp 049}, \runref{exp 051}): fresh Rademacher coins
  accumulated in exact ridge normal equations over $\phi$, with the unit-normalized frozen
  prior; per-state counts by statistics, valid under any visitation.
\end{enumerate}
Both are one solve away from ordinary deep RL practice (recursive least squares on the
penultimate features --- the same object the elliptical-bonus family already maintains), and
both inherit the trunk's generalization: the bonus is a function of the features, not of a
position table.

% >>> AUTO-GENERATED TABLE START: neural
\begin{table}[H]
\centering
\footnotesize
\setlength{\tabcolsep}{4pt}
\begin{tabular}{@{}>{\raggedright\arraybackslash}p{6.4cm} r r r r@{}}
\toprule
method & dev\_worst $\downarrow$ & dev\_mean & slope & slope std \\
\midrule
\multicolumn{5}{@{}l}{\textbf{maze + AntMaze states (uniform)} (10 seeds)} \\
deep trunk + exact shrink head & \textbf{0.038} & \underline{0.038} & -0.500 & 0.000 \\
deep trunk + exact coin-flip head & \underline{0.054} & \textbf{0.030} & -0.503 & 0.004 \\
gradient coin-flip net, Adam $10^{-2}$ (best gradient variant) & 1.148 & 1.159 & -0.696 & 0.264 \\
gradient coin-flip net, Adam $10^{-4}$ (paper rate) & 2.202 & 1.849 & -0.648 & 0.347 \\
\midrule
\multicolumn{5}{@{}l}{\textbf{held-out 20\% (vectors)} (10 seeds)} \\
deep trunk + exact coin-flip head & \textbf{1.678} & \textbf{0.264} & -0.609 & 0.549 \\
deep trunk + exact shrink head & \underline{1.935} & \underline{0.292} & -0.944 & 1.282 \\
\midrule
\multicolumn{5}{@{}l}{\textbf{nonuniform visitation (vectors)} (10 seeds)} \\
deep trunk + exact coin-flip head & \textbf{0.618} & \textbf{0.138} & -0.692 & 0.198 \\
deep trunk + exact shrink head & \underline{1.216} & \underline{0.711} & -0.352 & 0.145 \\
\bottomrule
\end{tabular}
\caption{The phase-2 neural ladder. Blocks are domain--regime combinations; within each block rows are sorted by dev\_worst (lower is better; best and second-best dev\_worst and dev\_mean bold and underlined). Every method starts at exactly 1 (start\_dev 0, omitted). In the held-out blocks the metric scores trained states against their counts AND the never-trained 20\% against the constant oracle value 1, so it punishes a bonus that generalizes the decay onto states never actually visited.}
\label{tab:neural}
\end{table}
% >>> AUTO-GENERATED TABLE END: neural

\subsection{What the generalization regimes add}
\label{sec:neural-generalization}

% (numbers spliced by the tick after the heldout/nonuniform experiments complete)

\clearpage

%==============================================================================
\section{Validation runs}
\label{sec:validation}

Every finalist and baseline was re-run at $30$ seeds on the three fixed point sets (the
top-right goal cell joins the two iteration sets), in both regimes; the dense grid validates
the champions' capacity behavior. Recall the deviation of \Cref{eq:dev}; \Cref{tab:validation}
reports the same metric battery as the ledger, now at validation seed counts.

% >>> AUTO-GENERATED TABLE START: validation
\begin{table}[H]
\centering
\footnotesize
\setlength{\tabcolsep}{4pt}
\begin{tabular}{@{}>{\raggedright\arraybackslash}p{5.6cm} r r r r r@{}}
\toprule
method & dev\_worst $\downarrow$ & dev\_mean & start\_dev & slope & slope std \\
\midrule
\multicolumn{6}{@{}l}{\textbf{uniform} (30 seeds)} \\
Hadamard coins (collision-free spikes), adaptive dictionary & \textbf{0.033} & \textbf{0.003} & 0.00 & -0.500 & 0.000 \\
residual-encoded shrink (exact solve) & \underline{0.038} & 0.038 & 0.00 & -0.500 & 0.000 \\
coin flips, adaptive dictionary & 0.054 & 0.030 & 0.00 & -0.502 & 0.003 \\
elliptical posterior readout ($\sigma \le 0.35$) & 0.139 & 0.046 & 0.00 & -0.499 & 0.001 \\
Hadamard coins (column-permuted), adaptive dictionary & 0.314 & \underline{0.008} & 0.00 & -0.501 & 0.001 \\
Adam $10^{-3}$, raw readout (control) & 1.653 & 1.250 & 1.62 & -1.038 & 0.356 \\
AdaGrad $3{\times}10^{-3}$ + initial-copy readout & 1.772 & 0.591 & 0.00 & -0.882 & 0.502 \\
SGD-$1/t$ (parent project's best), raw readout & 2.410 & 1.741 & 1.62 & -1.413 & 1.031 \\
\midrule
\multicolumn{6}{@{}l}{\textbf{nonuniform} (30 seeds)} \\
coin flips, adaptive dictionary & \textbf{0.466} & \underline{0.149} & 0.00 & -0.623 & 0.086 \\
Hadamard coins (column-permuted), adaptive dictionary & \underline{0.468} & \textbf{0.145} & 0.00 & -0.623 & 0.085 \\
elliptical posterior readout ($\sigma \le 0.35$) & 0.505 & 0.246 & 0.00 & -0.567 & 0.060 \\
residual-encoded shrink (exact solve) & 1.441 & 0.804 & 0.00 & -0.500 & 0.000 \\
AdaGrad $3{\times}10^{-3}$ + initial-copy readout & 1.857 & 1.031 & 0.00 & -0.898 & 0.487 \\
\midrule
\multicolumn{6}{@{}l}{\textbf{dense grid} (30 seeds)} \\
Hadamard coins (collision-free spikes), adaptive dictionary & \textbf{0.054} & \textbf{0.003} & 0.00 & -0.500 & 0.000 \\
Hadamard coins (column-permuted), adaptive dictionary & \underline{1.029} & \underline{0.008} & 0.00 & -0.501 & 0.001 \\
Hadamard coins, sign-only scramble (broken variant, kept as record) & 4.690 & 2.394 & 0.00 & -0.550 & 0.023 \\
\bottomrule
\end{tabular}
\caption{Validation at $30$ seeds. Uniform and nonuniform blocks pool the three fixed point sets; the dense-grid block is the $432$-point capacity stress. Within each block rows are sorted by dev\_worst (the arrowed sort column; lower is better), and the best and second-best dev\_worst and dev\_mean are bold and underlined. start\_dev is $0$ for every method with a normalized readout; the two raw-readout controls carry the untrained bonus spread. Slope columns are unmarked (closeness to $-1/2$, not size, is the goal).}
\label{tab:validation}
\end{table}
% >>> AUTO-GENERATED TABLE END: validation

Three validation-only observations:
\begin{enumerate}[label=(\roman*),leftmargin=*]
\item \textbf{The third point set breaks the optimizer family further.} AdaGrad's pooled slope
  spread grows from $0.20$ (two sets) to $0.50$ (three sets): the top-right goal cell ---
  already the heterogeneity flag of the parent project's convergence runs --- decays very
  differently under parameter-space optimizers. The exact constructions do not move
  ($0.0381$; $0.054$).
\item \textbf{The Hadamard refinement needed two bugs found, then set the uniform record.}
  Orthogonal coin blocks scrambled by signs alone send every full block's sum onto one
  coordinate, and consecutive blocks cancel: past $512$ visits the running-mean norm
  collapses (val 109--111, kept as the record of the broken variant). A per-block column
  permutation restores exactness per block, but two blocks of one position can still land
  their sums on the same coordinate ($\approx 5\%$ of positions over eight blocks), shifting
  that position's level (val 114--116). Drawing each position's block permutations with
  distinct spike coordinates removes that too, and the result is the uniform-regime record:
  dev\_worst $0.0335$, dev\_mean $0.0031$ --- below the shrink construction's off-by-one
  floor, because a Rademacher first draw has norm $\sqrt{d}$ exactly and orthogonal blocks
  make every later checkpoint exact as well (val 117--118).
\item \textbf{The nonuniform champion is statistical.} Coin-flip counting on the adaptive
  dictionary holds its iteration-phase score at $30$ seeds; the closed-form elliptical
  readout with the tighter bandwidth is second, and carries zero estimator noise in exchange
  for the overlap-count inflation.
\end{enumerate}

\subsection{Ablations and horizon stress}
\label{sec:ablations}

Five single-knob probes bound the champions' sensitivities, each at $10$ seeds
(\Cref{tab:ledger}'s protocol):
\begin{enumerate}[label=(\roman*),leftmargin=*]
\item \textbf{Coin dimension} ($d = 128$ against $512$): dev\_mean doubles ($0.060$ against
  $0.030$) --- the chi floor scales as $\sqrt{1/(2d)}$, exactly as computed in
  \Cref{sec:metrics}.
\item \textbf{Dictionary insertion radius} ($0.2$ against $0.05$): a dictionary coarser than
  the finest point spacing cannot separate the center square's positions and the score
  collapses ($1.06$ uniform, $0.99$ nonuniform). The insertion radius must sit below the
  finest state spacing the bonus should resolve.
\item \textbf{Elliptical bandwidth} ($\sigma \le 0.2 / 0.35 / 0.75$): dev\_worst improves
  monotonically toward orthogonality ($0.077 / 0.114 / 0.202$) --- the overlap-count
  inflation is the whole remaining error of the closed-form readout.
\item \textbf{Horizon} ($32{,}768$ steps, $8\times$ the protocol): the residual-encoded
  shrink holds slope $-0.500000$ with deviation $0.033$ over four and a half decades;
  AdaGrad drifts to slope $-0.63$ with deviation $1.30$ --- the optimizer family's
  calibration does not survive longer horizons.
\end{enumerate}

%==============================================================================
\section{Findings}
\label{sec:findings}

\begin{enumerate}[leftmargin=*]
\item \textbf{The requirement splits into level and rate, and the level is cheap.} A frozen
  copy of the initial predictor as a per-position denominator (or a zero head under a
  unit-normalized target or prior) starts every curve at exactly $1$ online. Nothing about
  requirement (1) needs the optimizer.
\item \textbf{No optimizer or global schedule can satisfy requirement (2).} Under any
  learning-rate schedule the tangent-kernel modes' log-decays keep the fixed ratios
  $\lambda_j/\lambda_k$ (\Cref{sec:mechanism-one}); the parent project's aggregate $-0.503$
  was a calibration of the spectrum's bulk, with per-position slopes from $-0.4$ to steeper
  than $-4$ around it. AdaGrad comes closest ($\pm 0.1$ on two point sets) and still breaks
  on the third.
\item \textbf{The uniform regime is solved by a residual-encoded functional step.} The
  per-visit map $r \mapsto r/\sqrt{1+r^2}$ --- whose state is the bonus value itself ---
  realized exactly by a minimum-change interpolation solve, puts every position on
  $(m+1)^{-1/2}$ simultaneously: slope $-0.500000$ at every position, deviation $0.0381$
  $=$ the off-by-one floor, self-correcting under capacity stress, at $10$ and at $30$
  seeds. Its inner solve is the load-bearing part: gradient-descent inner fitting fails on
  slow kernel modes, and inner Adam overshoots.
\item \textbf{Per-position counts under arbitrary visitation need statistics, not
  recurrences.} The shrink map's state is corrupted by feature interference under partial
  visitation and integrates the damage; the coin-flip running mean (fresh $\pm1$ targets per
  visit, exact least squares, adaptive dictionary) averages interference out and tracks each
  position's own count: dev\_worst $0.47$ against per-position targets where every
  fixed-schedule method scores $1.4$ or worse.
\item \textbf{Orthogonal coin sequences remove the statistical floor --- with a caveat.}
  Scrambled-Hadamard visit sequences make a position's first $n$ coins sum to squared norm
  exactly $nd$, deleting the chi fluctuation; the sign-only scramble fails catastrophically
  past one block and needs the per-block column permutation.
\item \textbf{Representation quality is a hard constraint, not a detail.} The project's
  zero-bias ReLU features are rank-deficient on the maze grid (rank $106/108$, condition
  $10^{10}$) and break every least-squares construction; a fixed radial-basis bandwidth
  cannot serve two spatial scales at once; the growing data-dependent dictionary (insert
  far-enough visited states, bandwidth from the local spacing) serves both scales and the
  $432$-point stress set.
\item \textbf{The user-suggested landscape lens does not bind here.} The loss-landscape
  literature's chaos-with-depth findings (and the skip-connection remedy) concern deep
  plain networks; a $4 \to 256 \to 128$ network is already in the shallow-wide regime those
  papers call benign, and the measured pathology is the kernel spectrum, not non-convexity
  (\Cref{sec:literature}).
\end{enumerate}

%==============================================================================
\section{Literature}
\label{sec:literature}

Full per-paper notes live in the campaign's \texttt{literature/} folder (six commissioned
surveys); this section keeps what the campaign used.

\paragraph{Why the decay rate is position-dependent.}
The linearized analysis behind \Cref{sec:mechanism-one} is the neural-tangent-kernel
description of gradient descent \citep{jacot2018neural_neurips}, whose per-eigenmode
convergence rates were made explicit for overparameterized two-layer networks by
\citet{arora2019fine_icml}, with global-convergence guarantees for exactly our
shallow-wide regime in \citet{du2019gradient_iclr}. The spectral bias phenomenon --- low
kernel modes learn first --- is documented empirically in \citet{rahaman2019spectral_icml}
and quantified per frequency in \citet{basri2019convergence_neurips}; the extension to
non-uniform data density \citep{basri2020frequency_icml} proves the local convergence time
scales with one over the local density, which is why plain gradient methods are already
mildly count-sensitive under non-uniform visitation (\runref{exp 012}). The radial pattern
in our diagnostics is the arc-cosine kernel's growing diagonal
\citep{cho2009kernel_neurips}. Fourier input features reshape the composed kernel toward a
stationary one \citep{tancik2020fourier_neurips} --- the delocalization route to uniformity
that our sphere-projection attempt (\runref{exp 008}) approximated poorly.

\paragraph{Schedules, averaging, and adaptive optimizers.}
The $1/t$-schedule and iterate-averaging theory the parent project leaned on is the
stochastic-approximation line of \citet{polyak1992acceleration}, with the non-asymptotic
constant-step least-squares rate of \citet{bach2013non_neurips} and the tail-averaging
refinement of \citet{jain2017parallelizing_jmlr}. All of these give $O(1/n)$ \emph{expected
squared error under gradient noise}; our base regime is deterministic full-batch descent,
where the averaged iterate decays as $1/n$ per mode only past each mode's own transient
(\runref{exp 004}). AdaGrad's cumulative accumulator \citep{duchi2011adaptive} is the
count-like member of the adaptive family, and Adam's forgetting \citep{kingma2014adam} is
why the canonical RND bonus saturates.

\paragraph{Count-style bonus constructions.}
The coin-flip construction the campaign made exact is the pseudo-count estimator of
\citet{lobel2023flipping_icml}; the pseudo-count idea originates with density models
\citep{bellemare2016unifying_neurips, ostrovski2017count_icml}, and the closed-form
elliptical readout is the episodic elliptical bonus family \citep{henaff2022exploration_neurips}
built on the linear-bandit confidence width \citep{abbasi2011online}. The bridge the
campaign exploits --- fitting random targets estimates a posterior width that a ridge
accumulator computes exactly --- is the prior-network posterior connection of
\citet{ciosek2019conservative}, with randomized prior functions
\citep{osband2018randomized_neurips} as the never-visited baseline mechanism. Two recent
RND descendants document the exact defects this campaign scores: distributional RND names
the unequal initial bonus (``initial bonus inconsistency'')
\citep{yang2024exploration}, and random-distribution distillation injects fresh target
noise to convert optimization decay into statistical decay \citep{fang2025exploration} ---
mechanism two in this document's terms.

\paragraph{Making per-point rates uniform on purpose.}
Natural-gradient/Gauss-Newton steps contract every residual mode at the same rate
\citep{zhang2019fast_neurips}; the campaign's minimum-change interpolation update is the
exact-solve limit of that idea with the schedule replaced by a residual-encoded map. The
only prior work we found whose stated goal is equalizing per-point residual decay is the
physics-informed-network line, where per-point loss weights are tuned to balance decay
rates \citep{chen2025selfadaptive_journal_of_computational_physics}; our inverse-error
reweighting probe (\runref{exp 006}) is a crude member of that family and inherited its
limitation (diagonal preconditioning cannot fix off-diagonal coupling).

\paragraph{The landscape question the user raised.}
\citet{li2018visualizing_neurips} shows loss-landscape chaos growing with \emph{depth} in
plain networks and disappearing with skip connections or width. A $4 \to 256 \to 128$
network sits in the shallow-wide regime that the paper itself, and the two-layer
convergence theory \citep{du2019gradient_iclr}, classify as benign --- so residual
connections are not the lever for this problem; the kernel spectrum is
(\Cref{sec:mechanism-one}). The original bonus construction remains
\citet{burda2019exploration}.

%==============================================================================
\section{What this document does not settle}
\label{sec:limitations}

\begin{enumerate}[leftmargin=*]
\item \textbf{Off-support behavior --- measured, and the metric does not score it.} The
  campaign metric evaluates on the visited point set; \Cref{fig:offsupport} shows what the
  three champions read \emph{between} the trained positions after $256$ uniform visits. The
  residual-encoded shrink produces a flat field (off-support mean $0.0625$ against
  on-support $0.0624$): its global features generalize the decay everywhere, so an exploring
  agent would find nothing marked novel anywhere. The coin-flip dictionary keeps elevated
  bonus between cells (off-support mean $0.096$, maximum $0.71$). The unit-normalized
  elliptical readout also flattens (maximum $0.0624$): normalizing the features erases the
  far-from-every-center signal that would otherwise mark unvisited space. A bonus that is
  perfect \emph{on} the support and blind \emph{off} it is not yet an exploration bonus;
  scoring the off-support field is the natural next metric.

\begin{figure}[H]
\centering
\includegraphics[width=\linewidth]{figures/offsupport_fields.pdf}
\caption{The bonus field on a quarter-cell grid after $256$ uniform full-batch visits to the
$108$ cell midpoints (white dots), one panel per champion. The visited-cell level is
$(257)^{-1/2} \approx 0.062$ everywhere; the panels differ in what they claim about the
space in between.}
\label{fig:offsupport}
\end{figure}
\item \textbf{The full reinforcement-learning loop.} The campaign isolates the bonus from
  policy feedback on purpose. Inside a SAC or PPO loop the visitation law is itself driven
  by the bonus; whether the count-exact methods improve exploration end to end is the
  parent project's next train-run question.
\item \textbf{Streaming state spaces.} The adaptive dictionary was stressed to $432$
  distinct states with a $1024$-center cap; continuous high-volume state streams need
  eviction or merging policies the campaign did not test.
\item \textbf{Pure-gradient deep-network versions.} The exact constructions use a linear
  head and a solve. The best gradient-only approximations measured here (AdaGrad at a tuned
  rate; the quartic loss) are honest but loose; closing that gap inside a standard deep-RL
  training stack is open.
\end{enumerate}

\bibliographystyle{plainnat}
\bibliography{bibliography}

\end{document}
